\section{Birational constructions and reduced-boundary models}
\label{bir:section}

This section supplies the ordinary dlt constructions needed in the
algebraic-dimension-zero argument.  We prove special termination through
dimension four, using arbitrary klt termination through dimension three,
and then apply Assumption~\ref{hyp:mmp} only to ordinary effective klt
fourfold pairs.  In particular, arbitrary dlt fourfold termination is not
an input.

\subsection{Conventions, comparison, and integral descent}
\label{bir:conventions}

All pairs in this section are ordinary rational pairs with effective
boundary and rationally invertible adjoint.  Auxiliary programs start
on normal globally strongly $\Q$-factorial compact K\"ahler spaces.
Thus every global coherent rank-one reflexive sheaf has an invertible
reflexive power.  This condition will be used for global sheaves; no
claim of $\Q$-factoriality on arbitrary analytic open subsets is made.
Divisors over a space mean divisorial places represented on proper
modifications, identified on common higher models.  This convention
does not distinguish places by their action on global meromorphic
functions.

\begin{definition}[Elementary birational steps]
\label{bir:elementary}
An elementary step for an adjoint $J=K_T+B$ is either a nontrivial
projective bimeromorphic divisorial contraction, with the pair pushed
forward, or a diagram
\begin{equation}
 T\xrightarrow{f} Z\xleftarrow{f^+}T^+,
 \label{bir:small-diagram}
\end{equation}
whose morphisms are projective, bimeromorphic, and small, with strict
transform boundary on $T^+$.  The base is normal compact K\"ahler, and
the source and next space are normal globally strongly $\Q$-factorial
compact K\"ahler spaces.  All curves contracted by the negative
morphism span one nonzero ray for degrees of global rational line
bundles.  The line $-J$ is relatively ample; in a small step $J^+$ is
relatively ample.  No positive-side Bott--Chern rank condition is
imposed in this definition.
\end{definition}

A projective morphism here has a relatively ample holomorphic line
bundle.  In the relative constructions, ``relatively nef'' means
nonnegative degree on curves in its fibers.  Absolute nefness continues
to mean analytic nefness.

We use common projective resolutions and the natural meromorphic
comparisons of canonical bundles, defined by Jacobians in local
canonical frames, together with boundary pullback.  No global
meromorphic canonical frame is presumed.  A global rank-one reflexive
sheaf is transported through a bimeromorphic correspondence by
pullback to a common resolution, proper direct image, and double dual.
For a small correspondence this is the unique reflexive strict
transform.

\begin{lemma}[Exceptional comparison of global sheaves]
\label{bir:sheaf-comparison}
Let $h:Y\to T$ be a projective modification of normal spaces and let
$\mathcal P$ be a global rank-one reflexive sheaf on $Y$.  Suppose that
$(h_*\mathcal P)^{**}$ is rationally invertible.  Then, for some positive
integer $l$, an invertible sheaf $M$ on $T$ and an integral
$h$-exceptional Weil divisor $E$ satisfy
\begin{equation}
 \mathcal P^{[l]}\simeq
 \bigl(h^*M\otimes\OO_Y(E)\bigr)^{**}.
 \label{bir:exceptional-sheaf}
\end{equation}
\end{lemma}

\begin{proof}
Choose $l$ so that the corresponding reflexive power of
$(h_*\mathcal P)^{**}$ is invertible.  This is
$(h_*\mathcal P^{[l]})^{**}$, since the two sheaves agree where $h$ is an
isomorphism, including the general points of all prime divisors of
$T$.  Call this line $M$.  Evaluation of local sections of
$h_*\mathcal P^{[l]}$, and their images in $M$, gives meromorphic
comparisons with $h^*M$.  These comparisons are the same on their common
domain and hence define a divisorial comparison globally.  Its orders
vanish away from the exceptional primes.  The resulting integral
exceptional divisor gives~\eqref{bir:exceptional-sheaf} in codimension
one and therefore everywhere by reflexivity.  Local meromorphic
generators are enough for this argument; it does not represent an
arbitrary global line bundle by a global divisor.
\end{proof}

We shall use the projective analytic negativity lemma
\cite{FujinoAnalyticMMP,BirationalGeometry}: if a rational Cartier
divisor on a projective bimeromorphic morphism is relatively nef and
has nonpositive pushforward, it is nonpositive.  The local-over-target
proof of \cite[Lemma~3.39]{BirationalGeometry} applies in this setting.

\begin{lemma}[Comparison in an elementary step]
\label{bir:comparison}
On a common smooth projective resolution of an elementary step, with
projections $p$ to the negative model and $q$ to the next model, the
natural adjoint comparison is
\begin{equation}
 p^*J=q^*J'+F,\qquad F\geq0,\qquad q_*F=0.
 \label{bir:effective-comparison}
\end{equation}
Log discrepancies do not decrease.  They increase strictly for a
place whose center on either side is contained over the non-isomorphism
locus in the contraction base.  Consequently an elementary step
preserves klt, respectively dlt, singularities.
\end{lemma}

\begin{proof}
The codimension-one comparison gives $q_*F=0$.  The signs of the two
adjoints make $-F$ nef over the contraction base, and hence over the
target of $q$.  Negativity gives $F\geq0$.

In a small diagram the non-isomorphism sets in $Z$ are the same on
both sides.  Otherwise, over a region where one side is an isomorphism,
smallness identifies the other adjoint with a pullback, contradicting
its relative ample sign on a contracted curve.  A non-isomorphism
fiber is positive-dimensional by normality.  Over any such point,
lift a negative curve to the common resolution.  Its $F$-degree is
strictly negative, so the fiber meets $\Supp F$.  A connected
projective fiber meeting the support of an effective relatively
anti-nef Cartier divisor is contained in that support: if a component
outside the support met it, a curve section through an intersection
point would have positive intersection.  Clearing denominators gives
the same statement for $F$.  Thus $\Supp F$ contains the full fibers
in question.  Pullback to any higher model now proves strict increase
at each specified place.

For dlt preservation use the characterization by log canonicity and an
SNC open set meeting every lc center.  A new discrepancy-zero place
was already a zero place, and strictness keeps the general point of
its center out of the surgery.  The SNC characterization, checked on
global log resolutions, is therefore preserved.  The klt assertion is
immediate.  Finally, on a globally strongly $\Q$-factorial dlt pair,
slightly decreasing the coefficients of the floor gives a klt pair:
the floor is effective rational Cartier, and every discrepancy-zero
place has center in it and positive order on its pullback.
\end{proof}

We recall several analytic projective facts used in these arguments.
Projective modifications and projective spaces over compact K\"ahler
bases are K\"ahler, also for singular spaces.  To see the needed
positivity directly, give a relatively ample line a metric by local
relative embeddings and Fubini--Study metrics on a finite base cover,
scaling back from the relatively very ample powers.  Patch weights on
the actual line by a partition from the base.  In a local embedding,
the active coordinates and frame changes lift to ambient holomorphic
germs.  Base cutoffs have zero first and second derivatives in
vertical directions, so the patched weight has positive Levi form on
vertical Zariski tangent vectors.  Add a large multiple of a pulled-back
base K\"ahler potential.  One multiplier works on compact regions:
otherwise a sequence of unit tangent vectors on shrinking compact
charts would limit to a vertical tangent vector contradicting the
strict vertical positivity.  The Zariski tangent spaces vary in a
closed set in a fixed local embedding.  Squared absolute values of
defining equations adjust the extensions to strict ambient
positivity.  This also treats finite morphisms, for which the trivial
line is relatively ample; compare \cite{VarouchasKahler}.

Common resolutions in projective diagrams are obtained by graph or
fiber products followed by projective log resolution.  Relative
ampleness composes after adding sufficiently large pullback multiples
from below, locally over compacta.  A curve on the base of a
projective surjection has a curve lift: pull back to its normalization,
a projective curve, and use relative generation and base twists to
make the total space projective; then take curve sections.  This
justifies all curve lifts above and below.  Individual projective
fibers permit the usual curve sections and the projective Kleiman
criterion.

The established local inputs are the relative klt cone and base point
free theorems for projective analytic morphisms over Stein
neighborhoods of compacta satisfying property (P), and the big-klt
adjoint finite-generation theorem.  We can take arbitrarily small
compact coordinate neighborhoods cut out by balls or polydiscs, with
Stein ambient neighborhoods and the required finite-component
intersection property.  We use the cone theorem with its finite
negative-ray decomposition after a positive relatively ample
truncation.  We use base point freeness in the following integral
form: if $L$ is relatively nef Cartier and
$aL-(K+\Gamma)$ is relatively ample for some positive integer $a$,
then every sufficiently high integral power of $L$ is relatively
generated after shrinking.  See
\cite[Theorems~6.2, 6.5, and 7.2]{FujinoAnalyticMMP}.

\begin{lemma}[Integral klt descent]
\label{bir:integral-descent}
Let $f:T\to Z$ be projective bimeromorphic with normal target.  Suppose
an ordinary klt adjoint is $f$-antiample.  If an integral line bundle
$L$ has degree zero on every contracted curve, then $f_*L$ is a line
bundle and evaluation is an isomorphism
\begin{equation}
 f^*(f_*L)\simeq L.
 \label{bir:integral-pullback}
\end{equation}
\end{lemma}

\begin{proof}
The line $L$ is relatively nef, and the base point free hypothesis
holds by fiberwise ampleness.  Locally over a smaller base neighborhood,
all sufficiently high powers are generated.  The induced maps are
constant on each connected fiber, since their tautological lines have
degree zero on every curve there.  They factor through $Z$: the graph
projection is finite bimeromorphic onto the normal base.  The pulled-back
tautological lines descend two consecutive powers of $L$, and their
quotient descends $L$ itself.  The projection formula identifies the
descent with $f_*L$ and the pullback map with evaluation.  These
intrinsic identifications agree on overlaps.
\end{proof}

\subsection{Finite generation, nonextraction, and continuation}
\label{bir:finite-generation}

The finite-generation result we use is
\cite[Theorem~3.1]{DHPFourfoldMMP}: on a smooth space projective over
the indicated analytic base, multigraded rings of rational adjoints
with simultaneous SNC effective subunit boundaries and a common
relatively ample rational part are locally finitely generated, after
clearing denominators.  The relevant sums are klt.  All applications
here are over bimeromorphic projective bases, so the relative bigness
conditions hold.  The analytic projective big-klt framework and
finiteness of models are also described in
\cite[Theorem~E]{FujinoAnalyticMMP}.  Neither citation is used as an
arbitrary termination theorem over a nonprojective base.

We spell out the reduction to that theorem.  Over a Stein base in the
bimeromorphic case, the direct image of either sign of an actual line
bundle is a coherent sheaf of generic rank one.  Cartan generation
therefore supplies a nonzero section after shrinking.  Thus both signs
have local-over-base effective meromorphic representatives.  Local
canonical representatives can be obtained in the same way, or from
forms pulled back from general local projections downstairs.

For finitely many effective rational klt boundaries containing a
common positive relatively ample part, take a simultaneous projective
log resolution $a:S\to T$.  Add effective exceptional corrections to
their log pullbacks so that the resulting SNC boundaries
$\Gamma'_j$ are effective and have exceptional coefficients strictly
between zero and one.  Choose an effective exceptional divisor $E$
with $-E$ resolution-ample, by composing the exceptional antiample
choices for the successive blowups.  If $H$ is the common effective
ample part below, then $a^*H-\beta E$ is relatively ample for sufficiently
small $\beta>0$.  A common small multiple can be removed from each
$\Gamma'_j$ while preserving effectiveness and subunit coefficients:
the strict transforms contain the common part, and the exceptional
coefficients have positive margins.  Relatively ample summands may
also be represented by divided free general divisors after relative
generation, Stein sections, and analytic Bertini.  The smooth theorem
applies.  Effective exceptional corrections leave the cleared adjoint
rings unchanged by projection to a normal space.  The finitely many
actual bundle identifications can be powered and tensored
simultaneously, so they identify the multigraded rings multiplicatively.

In particular, a single rational klt adjoint has locally finitely
generated ring in this setting.  Add a sufficiently small positive
rational multiple of the sum of effective representatives of opposite
relatively ample integral bundles.  Their sum is actually linearly
equivalent to zero.  The addition supplies the common ample part and
preserves klt on a fixed resolution near the compactum.  This
replacement uses the projective bimeromorphic Stein setting
essentially.

\begin{lemma}[Relative Proj extracts no divisors]
\label{bir:proj-nonextraction}
Let $X_0$ be normal and projective bimeromorphic over a normal compact
base $T_0$.  Let $L$ be a rational line whose cleared relative ring is
locally finitely generated.  The normalized main component $Y$ of the
relative Proj is projective bimeromorphic over $T_0$;
$X_0\dashrightarrow Y$ extracts no divisors, and the reflexive trace
of $L$ on $Y$ is an actual relatively ample rational line.
\end{lemma}

\begin{proof}
Compactness permits a common divisible integer $m$ such that the
relative ring of $mL$ is generated in degree one.  Resolve its base
ideal and graph, and denote the maps to $X_0$ and $Y$ by $p$ and $q$.
The map to Proj lifts to the normalization.  There is an actual
moving/fixed decomposition
\begin{equation}
 p^*(mL)=q^*H+G,\qquad G\geq0,
 \label{bir:proj-decomposition}
\end{equation}
where $H$ is the relatively ample tautological line on $Y$.  Degree-one
generation and projection show that every relative section of
$p^*(kmL)$ has vanishing at least $kG$.

First, $G$ is $q$-exceptional.  If a component mapped to a prime on
$Y$, then $q$ would be an isomorphism at its general point.  For large
$k$, relative ample generation of
$q_*\OO(G)\otimes H^{\otimes k}$ gives a local-over-base section having
a pole at that prime relative to $H^{\otimes k}$.  Multiplying its
pullback by the section of $(k-1)G$ gives a section of $p^*(kmL)$ whose
vanishing is less than $kG$, a contradiction.

Second, every $p$-exceptional prime $P$ is $q$-exceptional.  Otherwise
relative generation of $q_*\OO(P)\otimes H^{\otimes k}$ gives a section
with a pole at the image prime.  Add $kG$ to obtain a section of
$p^*(kmL)+P$.  Since $p_*\OO(P)=\OO_{X_0}$ by normality, it comes from
the original ring.  As a section of the enlarged line it must vanish
once along $P$, whereas the chosen pole prevents that vanishing.
Here the coefficient of $G$ at $P$ is zero by the first part.  This is
again a contradiction.  Thus no divisor is extracted.  Taking the
trace of~\eqref{bir:proj-decomposition} gives $L_Y=H/m$ as actual
rational lines.
\end{proof}

\begin{lemma}[Finitely many marked models]
\label{bir:finite-models}
For a rational polytope of actual adjoints satisfying the preceding
local multigraded finite-generation hypotheses, only finitely many
marked normalized main relative Proj models occur at rational
parameters.
\end{lemma}

\begin{proof}
After a common Veronese and shrinking, choose finitely many homogeneous
generators of the multigraded ring.  Diagonal rings on rational degree
rays are finitely generated by the semigroup argument for monomial
degrees.  Their Proj charts can be taken with homogeneous monomials as
denominators.  The subsets of generators that occur as supports of
monomials on the specified positive degree ray form a finite pattern.
Localization at such a monomial inverts exactly its support, and the
multidegree-zero subring is the degree-zero localization of the
diagonal ring.  Intersections use unions of supports with the
canonical localization maps.  Thus parameters with the same pattern
have the same charts and gluing.  Taking the main component and
normalization preserves finiteness.

The base identification fixes the marking on the common bimeromorphic
open.  A finite collection of the smaller neighborhoods covers the
compact base.  If two marked models agree on each neighborhood, their
identifications over the base agree on the common dense open, hence
everywhere, and glue uniquely.  There are consequently only finitely
many global marked possibilities.
\end{proof}

\begin{proposition}[Relative continuation in the global strong category]
\label{bir:continuation}
Let $(T,B)$ be an effective rational klt or dlt pair satisfying the
conventions above, and suppose $T$ is projective bimeromorphic over a
normal compact K\"ahler space $V$.  If $K_T+B$ is not curve-nef over
$V$, there is an elementary negative step over $V$.  Its next space
remains projective over $V$, compact K\"ahler, and globally strongly
$\Q$-factorial, and the appropriate singularity type is preserved.
\end{proposition}

\begin{proof}
Use the finite-dimensional space of degrees of global rational line
bundles on curves contracted over $V$; finite dimensionality follows
from first Chern classes in finite-dimensional cohomology.  For a
fixed relatively ample $H$, the closed curve cone has a compact slice
of $H$-degree one.  Indeed, for every global line $L$, both
$kH+L$ and $kH-L$ are relatively ample for sufficiently large $k$, so
all coordinates on the normalized slice are bounded.

Choose a negative extremal ray.  For dlt input decrease the floor
coefficients rationally just enough to obtain a klt adjoint $J_0$
still negative on this ray; for klt input let $J_0=K_T+B$.  Cover the
base by interiors of finitely many Stein compacta as above.  Projecting
the local cone decompositions to global degrees shows that, for each
positive rational $\delta$, the normalized slice lies in the convex
hull of its compact part $J_0+\delta H\geq0$ and finitely many points
of actual contracted curves.  A projected remainder stays in the
global closed cone, and a remainder of zero $H$-degree contributes
zero.  The stated convex hull is compact.

Choose $\delta$ so that the selected ray is strictly in the truncated
negative region.  It is therefore represented by an actual curve.
Separating its point on the slice from the compact hull of the
remaining generators and remainder gives a supporting nef class
vanishing only on this ray.  In the open set of such supports one can
choose a rational global line $N$ in the rational annihilator of the
ray, so that a positive multiple of $N$ minus $J_0$ is relatively
ample.  Strict positivity on the compact slice and fiberwise
ampleness give this last assertion.  Local klt base point freeness on
the finite cover generates a common multiple of $N$.  Its section
morphism and Stein factorization give a projective bimeromorphic
contraction $f:T\to Z$ over $V$, with normal target and connected
fibers, contracting precisely the ray.  The line $-J_0$ is
$f$-ample by fiberwise ampleness.

If an exceptional prime $E$ exists, it is rational Cartier and
negativity gives $E$ strictly negative degree on the ray.  All
contracted curves lie in $E$; they cover the nontrivial fibers.  There
can be no second exceptional prime, since its negative ray degree
would force the same curves, and hence $E$, into it.  For a global
rank-one reflexive sheaf on $Z$, take its reflexive pullback to $T$.
It is rationally invertible.  Add a rational multiple of $E$ to kill
its ray degree and apply Lemma~\ref{bir:integral-descent} after
clearing denominators.  The descended rational line agrees with the
given sheaf off codimension two on $Z$ and hence everywhere by
reflexivity.  Thus $Z$ has the required global strong property.

If $f$ is small, apply the preceding finite generation over $Z$ to
$J_0$.  Lemma~\ref{bir:proj-nonextraction} gives a projective positive
model with no extracted divisors; it is therefore small over $Z$.
The transformed adjoint is relatively ample and rationally invertible.
For any global line $L$ on $T$, killing the ray degree and applying
integral descent gives an actual rational identity
\begin{equation}
 L=cJ_0+f^*M,\qquad c\in\Q.
 \label{bir:one-ray-identity}
\end{equation}
Smallness gives the corresponding identity on the positive model.
Any global rank-one reflexive sheaf there first transports to $T$,
where a power is such an $L$.  Transforming back and using
\eqref{bir:one-ray-identity} proves the global strong property on the
positive model.  The positive morphism cannot be an isomorphism,
since that would make $J_0$ a pullback from $Z$.

For dlt input the unperturbed adjoint has negative degree on the ray.
Its version of~\eqref{bir:one-ray-identity} has $c>0$, so its transform
has the positive sign as well.  Lemma~\ref{bir:comparison} preserves
the required singularities.  All resulting spaces are projective
bimeromorphic over $V$ and hence compact K\"ahler, so the construction
can continue whenever relative curve-nefness fails.
\end{proof}

\subsection{Two finiteness counts and transversal adjunction}
\label{bir:counts}

\begin{lemma}[Divisorial count]
\label{bir:cycle-count}
A sequence of bimeromorphic transformations extracting no divisors
between normal irreducible compact K\"ahler $n$-folds contracts
divisors only finitely often.
\end{lemma}

\begin{proof}
Count the dimension of the span of prime $(n-1)$-cycle classes in
$H_{2n-2}(-,\R)$.  This is a finite nonnegative integer, since compact
analytic spaces are triangulable and have finite-dimensional homology.
Nonextraction provides common isomorphic opens whose complement in
the target has dimension at most $n-2$.  The localization sequence
identifies the target's degree-$2n-2$ homology with the Borel--Moore
homology of that open.  Restriction from the source maps its prime-cycle
span onto the target's span by strict transforms, killing the classes
of lost primes.  Each lost prime has nonzero class by its positive
K\"ahler volume, so the count strictly drops.

For completeness, that volume detects a topological class also on a
normal singular space.  K\"ahler potentials with pluriharmonic
differences determine a class in $H^2(-,\R)$ through the real-part
sequence
$0\to\ii\R\to\OO\to\mathcal{PH}\to0$, with a fixed normalization.
If pluriharmonicity is initially known only on the regular locus,
pull the difference to a resolution.  Its smooth pullback is
pluriharmonic.  Near the compact fiber over a point, choose holomorphic
real-part primitives and adjust imaginary constants to make their
values on that fiber agree.  The real part is constant there, and
each irreducible fiber germ forces the holomorphic primitive to be
constant there.  The primitives therefore agree on overlaps near
the fiber.  A finite cover and smaller neighborhoods glue them on a
neighborhood of the whole fiber, and properness and normality descend
them.  Thus the original difference is locally a holomorphic real
part.  Resolving a compact prime cycle now evaluates the corresponding
class power as the strictly positive integral of the pulled-back
K\"ahler form.  Its fundamental class pushes to the stated cycle
class.
\end{proof}

\begin{lemma}[Strict low-discrepancy counts]
\label{bir:low-count}
For a compact ordinary rational klt pair there exists
$0<\epsilon\leq1$ such that every discrepancy is at least $\epsilon$
and only finitely many exceptional places have log discrepancy
strictly less than $1+\epsilon$.  In particular all exceptional places
of discrepancy at most one can be realized on a single projective
SNC resolution.  For a terminal pair with largest boundary
coefficient $b_0$ (zero for empty boundary), the exceptional places
of discrepancy strictly less than $2-b_0$ are finite.  For an lc
pair the first assertion has the same form when restricted to
centers not contained in its non-klt locus.
\end{lemma}

\begin{proof}
On an SNC resolution let the log-pullback coefficients be $d_j<1$,
put $w_j=1-d_j$, and take
$\epsilon=\min(1,\min_jw_j)$, with value one if the list is empty.
For a place still exceptional over this resolution, choose at the
general point of its center normal coordinates $x_1,\ldots,x_s$,
the first $b$ defining the boundary components through the center.
The top exterior Jacobian calculation gives
\begin{equation}
 a_E\ \geq\ \sum_{j=1}^{b}w_j\ord_E(x_j)
             +\sum_{j=b+1}^{s}\ord_E(x_j).
 \label{bir:jacobian-bound}
\end{equation}
At most one differential saves one order by normal differentiation
along the place.  This proves the lower bound.  A place in the
strict cutoff must have center a stratum: an additional normal
coordinate would give at least $1+\epsilon$ (or at least two if no
boundary is present).

Blow up that closed stratum.  Its new SNC weight is the sum of the
passing weights, and the next center lies in the new exceptional
component.  Until the place is divisorial, only strata of codimension
at least two can occur; each next weight increases by at least
$\epsilon$.  The fixed strict cutoff bounds the chain length, and
there are finitely many strata at each stage.  The finitely many
exceptionals already on the starting resolution complete the count.
All these blowups can be made globally and projectively.

For a terminal pair, exceptional log-pullback coefficients are
negative.  The same calculation uses the strict cutoff $2-b_0$;
non-stratum centers and centers involving such exceptional
components cannot contribute new places below it.  For the lc
version take the minimum of one and the strictly positive weights.
A center not contained in the non-klt locus meets no zero-weight
component generically, so the same argument applies.
\end{proof}

\begin{lemma}[Transversal surface calculation]
\label{bir:transversal}
Let a normal analytic space carry an ordinary rational Weil boundary
with rationally Cartier adjoint.  At an analytically general point
of a codimension-two irreducible locus, a general transversal surface
cut is normal and the crepant formula on a simultaneous log resolution
restricts to its exact surface crepant formula.  For a coefficient-one
prime, normalized divisorial adjunction is computed by normalized
curve adjunction on this cut.
\end{lemma}

\begin{proof}
Choose local parameters $t_1,\ldots,t_{n-2}$ along the locus from an
embedding.  Take a sufficiently general nearby common value, regular
on the smooth locus, upstairs on a projective log resolution, and on
all relevant smooth strata; discard images of nondominating strata.
The cut upstairs is a smooth surface, with the required SNC support,
and no component is exceptional.  Codimension-two bad loci are cut
to isolated points.

The surface downstairs has dimension two and is regular in
codimension one.  It is Cohen--Macaulay as well.  Indeed the
non-Cohen--Macaulay locus of
a normal analytic space has codimension at least three, by the local
depth/coherent Ext criterion and normality at height at most two.
Thus at the chosen general point the parameters are a regular
sequence.  The cut is generically reduced and Cohen--Macaulay, hence
reduced; Serre's criterion gives normality.  Read the restricted
boundary transversely as a cycle and take its closure.  Complete
intersection adjunction off the isolated bad set, followed by
reflexive extension, identifies the restricted rational adjoint with
$K_S+B_S$, using division by the same parameter volume form upstairs
and downstairs.  The restriction of the crepant formula has the
correct nonexceptional coefficients.  Any difference from the
surface crepant formula is exceptional and relatively numerically
zero, hence zero by both signs of negativity.  This also restricts
klt formulas, or dlt formulas from resolutions preserving an SNC
good open, with their stated discrepancies.

For a coefficient-one prime, first subtract its smooth strict
transform in the resolved formula, then take residue and push to
its normalization.  At general points over the chosen
codimension-two locus, the cut of that strict transform is a smooth
curve germ, finite and generically an isomorphism onto its branch of
the reduced surface curve.  It is therefore the curve normalization.
Transversality preserves coefficient orders.  Restricting first to
the surface and then taking residue gives the same result.  Only
the indicated sums need be rational Cartier on the normalization;
the individual canonical and boundary terms need not be.
\end{proof}

\subsection{Extraction of prescribed low places}
\label{bir:extraction-section}

\begin{proposition}[Low extraction]
\label{bir:low-extraction}
Let $(T,B)$ be an ordinary rational klt pair on a normal globally
strongly $\Q$-factorial compact K\"ahler space.  Any specified set of
exceptional places of log discrepancy at most one can be extracted,
and no other exceptional prime extracted, by a projective crepant
morphism $Y\to T$ with effective klt boundary, where $Y$ is normal,
compact K\"ahler, and globally strongly $\Q$-factorial.
\end{proposition}

\begin{proof}
The set is finite by Lemma~\ref{bir:low-count}.  Choose a projective
SNC resolution $h:R\to T$ carrying it.  Put an effective SNC klt
boundary $\Theta$ on $R$, retaining the strict boundary and the
crepant coefficients of the specified primes, and choosing every
other exceptional coefficient strictly above its crepant value.
With $J=K_T+B$ this gives
\begin{equation}
 D(0)=K_R+\Theta=h^*J+P,\qquad P\geq0,
 \label{bir:extraction-origin}
\end{equation}
where $P$ is supported exactly on the undesired exceptionals.

Fix an $h$-ample line $H_0$ and let $E_1,\ldots,E_N$ be all
exceptional primes of $R$.  On a small full-dimensional rational
polytope in formal coefficient space consider
\begin{equation}
 D(u)=D(0)+u_0H_0+\sum_{i=1}^{N}u_iE_i,
 \qquad u_0\geq0,\quad |u_i|\leq\eta u_0.
 \label{bir:extraction-polytope}
\end{equation}
Choose $\eta>0$ small enough that the perturbation is relatively ample
at every nonzero parameter, and then make the polytope small enough.
On the finitely many Stein neighborhoods choose effective
representatives of $H_0,-H_0$, and $\pm E_i$.  A small common positive
multiple of the effective sum representing $H_0-H_0\sim0$ supplies a
common ample part.  Taking the other coefficients sufficiently small
on a simultaneous resolution gives equivalent ordinary klt adjoints
at the vertices.  Lemmas~\ref{bir:proj-nonextraction}
and~\ref{bir:finite-models} give finitely many marked normal relatively
ample models over $T$ for all rational parameters.

For each model occurring arbitrarily near zero, take the affine spans
of its parameter sets in successively smaller punctured
neighborhoods.  These nested affine spaces eventually stabilize.
Discard models not accumulating at zero.  If every eventual span
were proper, finitely many proper affine subspaces would cover all
sufficiently small rational points in the full-dimensional cone,
which is impossible.  Thus one marked model $Y$ occurs arbitrarily
near zero with full eventual affine span.

Choose affinely independent rational parameters for this model.
Their actual ample adjoint traces solve a rational linear system for
the traces of $D(0),H_0,E_1,\ldots,E_N$ in the group of rank-one
reflexive sheaves tensored with $\Q$.  All these traces are therefore
actual rational line bundles.  A sequence of its parameters tending
to zero makes $D(0)_Y$ curve-nef over $T$.  Taking
\eqref{bir:extraction-origin} in codimension one gives the actual
identity $D(0)_Y=h_Y^*J+P_Y$, with $P_Y$ effective exceptional.
Negativity forces $P_Y=0$.

No specified prime is lost.  At a nonzero parameter defining $Y$,
the relative system of a divisible power of $D(u)$ contains the
system of the ample perturbation, multiplied by the section of the
corresponding power of $P$ and a base pullback.  It embeds relatively
at general points outside $\Supp P$.  Every specified prime has
generic point there, including a discrepancy-one prime whose
crepant boundary coefficient is zero.  It cannot be contracted.
Thus $Y\to T$ is crepant with precisely the prescribed exceptional
primes and effective klt boundary.

Each surviving exceptional prime is rational Cartier by the traces
of the $E_i$.  For an arbitrary global rank-one reflexive sheaf on
$Y$, its trace on $T$ has an invertible power by the global strong
hypothesis.  Lemma~\ref{bir:sheaf-comparison} expresses the
corresponding power upstairs as a pullback line with an integral
exceptional correction.  Clearing the rational Cartier denominators
of that correction makes a further power invertible.  This proves
the global strong property.  Projectivity over the compact K\"ahler
base gives the remaining category assertions.
\end{proof}

\subsection{Arbitrary klt birational termination through dimension three}
\label{bir:threefolds}

\begin{proposition}[Terminal threefold termination]
\label{bir:terminal-termination}
An arbitrary sequence of elementary birational steps for an effective
rational terminal pair of dimension at most three is finite.  Here
terminal means that all exceptional log discrepancies are greater
than one.  No pseudo-effectivity hypothesis is required.
\end{proposition}

\begin{proof}
In dimensions at most two there are no nontrivial small diagrams, and
Lemma~\ref{bir:cycle-count} handles divisorial contractions.  In
dimension three discard a finite prefix to make all steps small.
Terminality persists by Lemma~\ref{bir:comparison}.  The underlying
threefold singularities are ordinary terminal: deleting the effective
rational Cartier boundary does not decrease discrepancies.  In
particular they are smooth at general curve points.  This conclusion
also holds for local analytic places.  On a global resolution the
exceptional coefficients for the empty boundary are negative; local
places still exceptional over that smooth resolution have ordinary
log discrepancy at least two.

Let $b$ be the largest boundary coefficient, or zero for empty
boundary.  The coefficient set is fixed under small steps.  If $b>0$,
test each step whose positive side has a flipped curve contained in
a coefficient-$b$ component.  If $b=0$, test every step with any
flipped curve; a nontrivial positive side exists by the ample-sign
argument in Lemma~\ref{bir:comparison}.  Blowing up the generic point
of that curve on the positive side gives an exceptional place with
log discrepancy
\begin{equation}
 t=2-\sum_j m_jb_j\leq2-b,
 \label{bir:terminal-test}
\end{equation}
where $m_j$ are nonnegative integral multiplicities.  One can realize
this place by the normalized blowup of the whole curve.  Positivity
and the fixed boundary denominators make the possible $t$ a finite
set.  For each such $t$, the number of exceptional places with
discrepancy strictly below $t$ is finite by
Lemma~\ref{bir:low-count}.  These counts do not increase.  At a tested
step its tested place had discrepancy strictly below $t$ before the
step and equals $t$ afterwards, so the corresponding count drops.
Only finitely many tested steps occur.

If $b>0$, consider the normalizations of the finitely many
coefficient-$b$ surfaces on the remaining tail.  Both sides map
bimeromorphically to the normalization of their common image in the
contraction base.  The positive map contracts no curve and is
therefore an isomorphism.  The negative normalization consequently
maps projectively bimeromorphically to the next normalization,
contracting a curve whenever the corresponding surface contains a
flipping curve.  These are compact K\"ahler surfaces.  The cycle
count leaves a tail on which no maximal-coefficient component
contains a contracted curve on either side.

Remove all maximal-coefficient components from the boundary on this
tail.  The new adjoint is still negative, because an effective
rational Cartier divisor has nonnegative degree on a curve not
contained in it.  Its one-ray identity with the original adjoint,
obtained by Lemma~\ref{bir:integral-descent}, has a positive rational
coefficient.  Smallness then makes its transform relatively positive.
Terminality persists after decreasing the boundary.  Induct on the
number of distinct positive boundary coefficients.  The empty-boundary
case was handled by testing every step, so the sequence is finite.
\end{proof}

We require a uniform index statement to pass from klt to terminal
pairs.  The local terminal-point results used are the classical
complex analytic threefold theorems of Mori--Reid and Kawamata.  If a
terminal point has canonical index $r>1$, there is an exceptional
place centered there with ordinary log discrepancy $1+1/r$; see
\cite{KawamataTerminalDiscrepancy}.  Also, the index of every integral
$\Q$-Cartier Weil divisor germ divides $r$, including when $r=1$.
For the latter statement, the analytic index-one cover is smooth or
an isolated cDV hypersurface, with simply connected punctured small
neighborhood.  The connectivity theorem for isolated hypersurface
links and the analytic Kummer sequence make its germ divisor class
group torsion-free.  Pulling up a $\Q$-Cartier divisor therefore makes
it principal, and the norm gives the divisibility downstairs.  The
small-discrepancy place can be detected globally on a compact
terminal space: on a global resolution the empty-boundary exceptional
coefficients are negative, so a local place still exceptional over
the resolution has discrepancy at least two.  A place of discrepancy
$1+1/r<2$ must already be a component of its restricted exceptional
divisor and hence a global exceptional place.

\begin{lemma}[A surface estimate for a single extraction]
\label{bir:surface-estimate}
Let $h:U\to T$ extract only a prime $P$ from an effective rational
klt threefold pair as in Proposition~\ref{bir:low-extraction}.  Write
its crepant boundary as $B^{\mathrm{str}}+(1-a)P$, where $0<a\leq1$.
There is a contracted curve $C\subset P$, not contained in
$\Supp B^{\mathrm{str}}$, such that
\begin{equation}
 (K_U+P)\cdot C\geq-3.
 \label{bir:surface-degree}
\end{equation}
\end{lemma}

\begin{proof}
First $-P$ is $h$-ample.  Compare any $h$-ample line with its
rationally invertible trace downstairs.  Their difference is a
rational multiple of the unique exceptional prime $P$, by
Lemma~\ref{bir:sheaf-comparison}.  Negativity makes that multiple
strictly negative, proving the assertion.

Divisorial adjunction of $K_U+P$ to the normalization $P^\nu$ gives
an actual rational line of the form $K_{P^\nu}+\Delta$, with
$\Delta\geq0$, even though $(U,P)$ need not be lc.  We justify the
sign, without requiring the individual terms on $P^\nu$ to be
rational Cartier.  Take residue along the smooth strict prime on a
log resolution and push its restricted crepant boundary to $P^\nu$.
This identifies the actual restricted line in codimension one, then
by reflexivity.  Lemma~\ref{bir:transversal} computes a tested
coefficient on a klt normal surface germ with the reduced curve of
all sliced branches of $P$.  Such a surface germ is a quotient of a
smooth germ by a small finite group, by the analytic klt surface
classification; see \cite[Chapter~4]{BirationalGeometry}.

On the smooth chart, adjunction to a normalized branch of a reduced
plane curve has effective different: the conductor contribution
from normalization and the intersections with other branches are
nonnegative.  More explicitly, hypersurface adjunction makes the
plane curve dualizing sheaf a line, and finite duality identifies the
normalization's dualizing sheaf with its pullback multiplied by the
conductor ideal.  The quotient pullback of the reduced divisor is
reduced.  Frames of a Cartier log-adjoint power pull to the
corresponding frames by the codimension-one \mbox{\'etale} property,
and their meromorphic residues are related by pluricanonical
pullback.  Therefore the different downstairs pulls to the
different upstairs plus the ramification divisor on the normalized
branch.  This proves $\Delta\geq0$.

On a minimal smooth resolution $\pi:S\to P^\nu$ the restricted
adjoint is $K_S+\Delta_S$ with $\Delta_S\geq0$.  Indeed $K_S$ is
relatively nef by curve adjunction, negative definiteness, and the
absence of exceptional smooth rational $(-1)$-curves.  The correction
$\Delta_S=\pi^*(K_{P^\nu}+\Delta)-K_S$ is relatively anti-nef with
effective pushforward; negativity applied to $-\Delta_S$ gives the
sign.

If $P$ maps to a curve, take a general smooth fiber on $S$, after
Stein factorization.  Its $K_S$-degree is at least $-2$.  If $P$ maps
to a point, $P$ and $S$ are projective.  The classification of smooth
projective surfaces supplies moving curves covering $S$ with
$K_S$-degree at least $-3$: use ample curves if the minimal canonical
class is nef, ruling fibers in the ruled case, or lines on $\PP^2$,
and general strict transforms; see \cite{BHPV}.  In either case
choose a curve not contained in the correction or in the finitely
many excluded curves, including the inverse image of
$\Supp B^{\mathrm{str}}\cap P$.  It maps birationally to a curve $C$
and the effective correction has nonnegative degree there.  This
proves~\eqref{bir:surface-degree}.
\end{proof}

\begin{proposition}[Uniform global reflexive index]
\label{bir:index}
Fix $0<\epsilon\leq1$ and an integer $n\geq0$.  For effective
rational klt threefold pairs on normal globally strongly
$\Q$-factorial compact K\"ahler spaces, suppose every discrepancy is
at least $\epsilon$, at most $n$ exceptional places have discrepancy
at most one, and every exceptional discrepancy greater than one is
at least $1+\epsilon$.  There is an integer $I(n,\epsilon)>0$ such
that $\mathcal F^{[I(n,\epsilon)]}$ is invertible for every global
rank-one reflexive sheaf $\mathcal F$ on each such space.
\end{proposition}

\begin{proof}
For $n=0$ the underlying space is terminal with the ordinary
exceptional discrepancy gap $1+\epsilon$.  The terminal-point
statement above bounds every canonical index by
$\lfloor1/\epsilon\rfloor$.  A global rank-one reflexive sheaf has
$\Q$-Cartier germs by the global strong hypothesis, so its local
indices divide those canonical indices.  Their bounded common
multiple gives $I(0,\epsilon)$.

Induct on $n$, seeking a multiple of previous bounds.  If there is no
low place use the previous case.  Otherwise extract only a place
$P$ of discrepancy $a\in[\epsilon,1]$ by
Proposition~\ref{bir:low-extraction}.  The pair on $U$ has effective
crepant boundary $B^{\mathrm{str}}+(1-a)P$.  Its exceptional places
are exceptional downstairs, except that $P$ is no longer counted;
all hypotheses hold with $n-1$.  Set $I'=I(n-1,\epsilon)$.

For the curve in Lemma~\ref{bir:surface-estimate}, crepancy and
effectiveness of $B^{\mathrm{str}}$ give
\begin{equation}
 -3\leq (K_U+P+B^{\mathrm{str}})\cdot C
       =aP\cdot C<0.
 \label{bir:index-degree}
\end{equation}
Let $\mathcal F$ be a global rank-one reflexive sheaf on $T$, and
let $\mathcal F_U$ denote its reflexive sheaf pullback.  Distinguish
this from the rational line pullback defined by an invertible power
of $\mathcal F$.  Their meromorphic comparison has the actual
rational-line form
\begin{equation}
 h^*\mathcal F=\mathcal F_U+sP.
 \label{bir:index-correction}
\end{equation}
Both $I'\mathcal F_U$ and $I'P$ are integral lines.  Degree zero of
the left side on $C$ gives
\begin{equation}
 s=\frac{I'\mathcal F_U\cdot C}{-I'P\cdot C},
 \qquad 1\leq -I'P\cdot C\leq\frac{3I'}{\epsilon}.
 \label{bir:index-denominator}
\end{equation}
The numerator and denominator are integers.  Put
\begin{equation}
 I=I'\operatorname{lcm}
       \{1,\ldots,\lceil3I'/\epsilon\rceil\}.
 \label{bir:index-recursion}
\end{equation}
Then $I$ and $Is$ are multiples of $I'$, so
$I\mathcal F_U+(Is)P$ is an integral line of zero degree on all
contracted curves.  Increase the crepant coefficient of $P$ by a
sufficiently small positive rational number.  The pair remains klt,
and its adjoint is relatively antiample because $-P$ is ample.
Lemma~\ref{bir:integral-descent} descends this integral line.  Its
descent agrees with $\mathcal F^{[I]}$ away from the codimension-two
center of the extraction, hence everywhere by reflexivity.  This
proves the induction.
\end{proof}

\begin{theorem}[Klt termination through dimension three]
\label{bir:klt-termination}
Every arbitrary sequence of elementary birational steps for effective
rational klt pairs of dimension at most three is finite, without a
pseudo-effectivity hypothesis.
\end{theorem}

\begin{proof}
Only dimension three remains.  Suppose there is an infinite
sequence; after the cycle count all steps are small.  The finite sets
of exceptional places of discrepancy at most one are nonincreasing,
so stabilize to a set $\mathcal E$.  There is a uniform positive
$\epsilon$ of the type in Proposition~\ref{bir:index} on this tail.
Indeed start with Lemma~\ref{bir:low-count} on its first model.  Only
finitely many places lie below its strict cutoff $1+\epsilon_0$;
outside $\mathcal E$ their discrepancies are greater than one, with
a positive minimum gap.  Take the minimum of this gap,
$\epsilon_0$, and the positive starting discrepancy lower bound.
All subsequent comparisons are nondecreasing.

The boundary denominators are fixed.  Proposition~\ref{bir:index}
gives a common denominator for the actual adjoints on all models of
the tail, and hence for all discrepancies by the crepant formulas.
The nondecreasing values on $\mathcal E$, bounded above by one,
therefore eventually become constant.  Over the first model of this
stabilized tail extract exactly $\mathcal E$ crepantly.  The
resulting pair is terminal with effective boundary and is globally
strongly $\Q$-factorial.

Lift a flip $X_-\dashrightarrow X_+$ over $Z$ from its extracted
terminal model $Y_0\to X_-$.  Run relative elementary steps over $Z$
using Proposition~\ref{bir:continuation}.  On a common projective
resolution of any finite prefix ending at $Y$, denote the pulled-back
adjoints of $Y_0,Y,X_+$ by $P_0,P,P_+$ respectively.  The comparisons
give
\begin{equation}
 P_0=P+G=P_++F,
 \qquad G,F\geq0,
 \label{bir:terminal-lift}
\end{equation}
where $G$ is exceptional over $Y$ and $F$ is exceptional over $X_+$.
The line $P_+$ is nef over $Z$.  Over $Y$ the divisor
$G-F=P_+-P$ is nef and has nonpositive pushforward; negativity gives
$G\leq F$, before any termination is known.  A lost prime of $Y_0$
would have positive coefficient in $G$, but zero coefficient in $F$:
this follows from smallness below for old primes and from constancy
of discrepancies on $\mathcal E$ for the extracted ones.  Thus no
prime is lost.  The run stays small and terminal, so
Proposition~\ref{bir:terminal-termination} makes it finite.  Its
endpoint has $P$ nef over $Z$.

At the endpoint negativity over $X_+$ applied to $F-G=P-P_+$ gives
the opposite inequality, hence $P=P_+$.  Since the adjoint on $X_+$
is relatively ample, its projection from the common resolution is
constant on the fibers over $Y$: otherwise a curve in a connected
projective fiber would have positive pullback degree.  Factoring the
graph over the normal space $Y$ gives a projective crepant morphism
$Y\to X_+$.  It extracts the same places, by smallness below and the
absence of lost primes, so the construction repeats.

Each nontrivial flip below forces a negative step in its lift, by
lifting a negative curve.  Infinitely many flips would concatenate
to an infinite small terminal sequence with effective strictly
transformed boundary, contrary to
Proposition~\ref{bir:terminal-termination}.  This proves the theorem.
\end{proof}

\subsection{Adjunction and special termination for dlt pairs}
\label{bir:dlt-section}

\begin{lemma}[Actual adjunction on normalized strata]
\label{bir:stratum-adjunction}
Let $(V,A)$ be an ordinary effective rational dlt pair, and let $T$
be the normalization of an lc center, with map $\nu:T\to V$.  A
chain of coefficient-one primes through its generic SNC stratum
gives an effective dlt adjunction pair $(T,A_T)$ and an actual
rational-line identification
\begin{equation}
 J_T=K_T+A_T=\nu^*(K_V+A).
 \label{bir:stratum-line}
\end{equation}
Its lc centers map to proper lc subcenters of the given center.
If the ambient coefficients belong to a DCC subset of $[0,1]$, the
adjunction coefficients belong to a DCC set depending only on that
set and the chain length.  In sufficiently divisible even powers,
the identification is the canonical meromorphic residue
identification, independent of the resolution and of the order of
the generic residues.  No global strong $\Q$-factoriality of $T$ is
asserted.
\end{lemma}

\begin{proof}
The lc centers are the images of the finitely many coefficient-one
strata on a log resolution, and each is generically an SNC floor
stratum.  Order the floor components defining the chosen generic
chain.  Take a projective log resolution isomorphic over an SNC
open meeting all lc centers.  Successively restrict its crepant
formula to the smooth strict strata by residue, then push the
adjunction boundaries in codimension one to their normalizations.
At each stage the SNC log pullback has coefficients at most one,
and its rational adjoint is the restriction of the ambient actual
line.  Reflexive extension gives~\eqref{bir:stratum-line}; any
exceptional difference from the crepant comparison is numerically
zero and vanishes by both signs of negativity.  A discrepancy-zero
stratum of a restricted formula comes from an ambient
coefficient-one stratum whose image meets the good open.  This gives
the sub-dlt discrepancy and good-open properties, and identifies
nested lc centers by the generic SNC calculation.

There is also a canonical meromorphic comparison, not merely an
abstract equality of line classes.  On a smooth strict stratum the
residue of the ambient log-pullback frame has exactly the pole and
zero orders of the remaining crepant boundary.  It therefore gives
the frame of the pushed adjunction in codimension one on the
normalization.  Extending the line identification reflexively and
composing with the natural embedding of its divisorial adjoint
sheaf into meromorphic pluricanonical tensors gives the claimed
map.  Different resolutions give the same map on the dense generic
SNC stratum, and hence everywhere meromorphically.  Permuting the
residues changes only signs, removed in even degree.  Sequential
adjunction through normalized intermediate strata agrees with this
strict-chain calculation by the same dense-open test.  In
particular all codimension-one boundary orders agree.

We prove effectivity and the DCC assertion at one restriction step;
iteration then proves the statement.  By
Lemma~\ref{bir:transversal}, a tested different coefficient is a
normalized-curve coefficient on a normal dlt surface with a
coefficient-one branch.  If the tested point is an lc center, the
surface pair is SNC there.  Indeed on a sliced log resolution
preserving the good open, a zero center must lift to the intersection
of two strict floor curves, with no exceptional locus through it;
the resolution is an isomorphism there.

Otherwise deleting the boundary leaves a numerically klt surface
germ.  The Mumford numerical pullbacks of effective terms are
effective by negative definiteness.  We recall why this is ordinary
klt even if rational Cartierness of the separate canonical term has
not yet been established.  On the minimal resolution of the
singular germ, write
\[
 M=(E_i\cdot E_j)<0,\qquad
 K_S-\pi^*_{\mathrm{num}}K=\sum_i a_iE_i.
\]
The numerical klt condition gives $a_i>-1$.  Relative nefness of
$K_S$, from minimality and curve adjunction, gives $Ma\geq0$ and
hence $a_i\leq0$.  For any nonzero effective integral exceptional
cycle $Y=\sum n_iE_i$, set $c_i=n_i+a_i>0$ on its support and
$c_i=0$ elsewhere.  Off-diagonal entries of $M$ are nonnegative, so
\begin{equation}
 c^{\mathsf t}M(n+a)\leq c^{\mathsf t}Mc<0.
 \label{bir:numerical-surface}
\end{equation}
Thus some supported $E_i$ has $(K_S+Y)\cdot E_i<0$.
Consequently $\OO_{E_i}(-Y+E_i)$ has degree greater than
$2p_a(E_i)-2$ and has zero $H^1$ by curve duality.  Peeling off such
components with the cycle exact sequences proves
$H^1(\OO_Y)=0$ for every exceptional cycle $Y$.  Grauert formal
functions gives rationality; multiples of the full exceptional
curve are cofinal with the maximal-ideal thickenings.

An integral multiple of the numerical pullback of any Weil divisor
is now a line-bundle divisor on the resolution with all exceptional
degrees zero.  Its class restricts to zero in $H^2$ of the
exceptional curve, by normalization and the degree description of
the curve's top cohomology.  Continuity around the compact fiber,
or topological proper base change, makes that class zero after
shrinking.  On the preimage of a small Stein neighborhood,
$H^1(\OO)=0$ by rationality.  The exponential sequence makes this
line actually trivial there.  Pushing in codimension one proves
that the original Weil divisor is rational Cartier; in particular
the canonical divisor is.  The numerical klt test is therefore the
ordinary one.

The analytic log-terminal surface classification now realizes the
germ as a smooth germ modulo a small finite group
\cite[Chapter~4]{BirationalGeometry}.  Pull up the full boundary by
canonical pullback, which is \mbox{\'etale} in codimension one.  No
exceptional place over the origin on the chart has nonpositive
discrepancy, by finite discrepancy comparison with positive
ramification factor.  The point blowup therefore tests total
boundary multiplicity strictly below two.  There is exactly one
smooth coefficient-one branch.  The finite group preserves it and
acts faithfully on its tangent: finite actions linearize, and a
nonidentity element with trivial tangential eigenvalue would be a
quasi-reflection.  Thus the group is cyclic of order $r$, also the
branch ramification index; $r=1$ is allowed.

Residue and ramification give $1-1/r$ for the branch alone.  Other
components of coefficients $d_j$ contribute intersection
multiplicities $k_j\geq0$ on the chart.  The different coefficient
is consequently
\begin{equation}
 1-\frac1r+\frac{\sum_jk_jd_j}{r},
 \qquad 0\leq\sum_jk_jd_j\leq1.
 \label{bir:different-coefficients}
\end{equation}
The upper bound follows from the sub-lc resolution formula.  This
proves effectivity and hence the full dlt condition at this step.
Positive elements of a nonnegative DCC set have a positive minimum,
if any occur.  The sums in~\eqref{bir:different-coefficients} thus
have a bounded number of positive terms, counted with multiplicity,
and satisfy DCC.  In a decreasing sequence of coefficients below
one, $r$ is bounded as well, since $1-1/r$ is a lower bound.  This
proves DCC for each restriction and for its iterations.
\end{proof}

\begin{lemma}[Strict comparison on strata]
\label{bir:stratum-comparison}
Suppose a small dlt elementary step preserves the generic points of
an lc center and its chosen adjunction chain.  The normalized
strata on both sides map projectively bimeromorphically to the
normalization $S$ of their common image in the contraction base.
Their adjoints are relatively antiample and ample, respectively.
On a common higher model their canonical pullback comparison
satisfies
\begin{equation}
 P_T-P_{T^+}\geq0.
 \label{bir:stratum-increase}
\end{equation}
The resulting discrepancy increase is strictly positive at any
place whose center on either stratum maps into the ambient
exceptional locus on that side.
\end{lemma}

\begin{proof}
The morphisms are the restrictions of the ambient projective
morphisms followed through finite normalizations.  The relative
ample signs restrict as stated.  Choose a common ambient resolution
preserving the generic SNC chain.  By
Lemma~\ref{bir:comparison} its effective pullback difference has
support containing the full fibers over the non-isomorphism set.
Restrict the two crepant formulas along the common strict chain.
The same coefficient-one terms are subtracted, and
Lemma~\ref{bir:stratum-adjunction} identifies the remaining canonical
comparisons.  Their difference is precisely the restriction of that
effective divisor.  The stratum itself is not contained in its
support, since its generic point is preserved.  The full inverse
image of any indicated center is contained in the support, so
pullback gives strictly positive multiplicity at each such place.
This proves both assertions for actual crepant boundaries, rather
than only for restricted first Chern classes.
\end{proof}

\begin{theorem}[Dlt special termination and modifications]
\label{bir:special-termination}
The following hold for ordinary rational pairs in the analytic
projective setting above.
\begin{enumerate}[label=\textup{(\roman*)}]
\item In dimension $d\leq4$, any dlt elementary sequence has a tail
whose exceptional loci on both sides are disjoint from every lc
center.
\item An ordinary rational lc pair of dimension $d\leq4$ on a
normal compact K\"ahler space, with rationally invertible adjoint,
has a projective crepant dlt modification whose total space is
globally strongly $\Q$-factorial and compact K\"ahler.  The boundary
is effective, and every exceptional prime has coefficient one.
\item For $d\leq3$, every dlt elementary sequence is finite.
\end{enumerate}
The bases of the elementary steps may vary.  Assertion \textup{(iii)}
is not asserted in dimension four.
\end{theorem}

\begin{proof}
We induct simultaneously on dimension.  The order within a
dimension is (i), then (ii), then (iii) when $d\leq3$.  Dimension
zero is immediate.  In proving (i) in dimension $d$, discard a
prefix so that all steps are small, by
Lemma~\ref{bir:cycle-count}.  A zero-discrepancy place on a new
space was already a zero place.  Strictness in
Lemma~\ref{bir:comparison} excludes containment of its center in
the exceptional locus.  The finite lists of lc centers therefore
stabilize, and the generic point of every remaining center is
preserved at every subsequent step.

Induct increasingly on the dimension $e$ of a surviving center;
the point case is already settled by preservation of its generic
point.  Normalize its successive transforms and use the same
generic adjunction chain.  By the smaller-center conclusion, the
stratum diagrams and adjunction pairs are isomorphisms near their
non-klt loci on a tail.  Their discrepancies are nondecreasing by
Lemma~\ref{bir:stratum-comparison}.

If a prime is extracted by a stratum transformation, it lies on
the new side over the ambient exceptional locus.  Its earlier
discrepancy is strictly smaller than its new discrepancy, the latter
being at most one by effectivity.  At the start of this tail it
therefore had discrepancy less than one and center outside the
non-klt locus, since the diagrams are unchanged near that locus.
Lemma~\ref{bir:low-count} gives finitely many possible exceptional
places; the possible nonexceptional positive-boundary primes add
only finitely many.  For each such place, repeated extraction would
give a strictly decreasing sequence of boundary coefficients, by
the strict comparisons.  The fixed adjunction DCC set of
Lemma~\ref{bir:stratum-adjunction} excludes infinitely many such
occurrences.  After discarding a prefix there are no stratum
extractions.

The cycle count then leaves no prime contractions by stratum
transformations either; for normal curves the diagrams are already
isomorphisms.  The primes now match.  Their coefficients do not
increase, and finite positive support and DCC stabilize the whole
boundary.  Neither morphism of a stratum diagram to $S$ can
contract a prime on this tail: its generic point would lie over
the ambient exceptional locus, and strict discrepancy increase
would contradict the matched coefficients.  Thus the stratum
diagrams are small, allowing isomorphisms, and their effective
pullback differences are exceptional over the positive stratum.
If the stratum transformation is an isomorphism, its discrepancies
agree, and strictness excludes any ambient exceptional intersection
with the stratum.  One may test a divisor over a point of such an
intersection, or the point itself for a curve.

For each remaining nontrivial stratum diagram, its dimension is
$e<d$.  Use the modification assertion already proved in dimension
$e$ to start on a projective globally strong crepant dlt model of
the negative stratum.  Run relative elementary steps over $S$ by
Proposition~\ref{bir:continuation}.  The lower-dimensional
termination assertion makes this run finite, with a relatively nef
endpoint.  On a common higher model write
\begin{equation}
 P_0=P+G=P_++F,
 \label{bir:dlt-lift}
\end{equation}
where $P_0,P,P_+$ are the initial, endpoint, and positive-stratum
adjoint pullbacks.  The divisors $G,F$ are effective and exceptional
over the respective endpoint sides.  Both $P$ and $P_+$ are
relatively nef, so negativity in both directions gives $G=F$ and
$P=P_+$.  Relative ampleness on the positive stratum makes its
projection constant on the connected projective fibers over the
endpoint, as in the proof of Theorem~\ref{bir:klt-termination}.
The graph therefore factors into a projective crepant morphism to
that stratum.  Remaining exceptional primes came from old
exceptional primes, since there was no extraction upstairs and the
stratum transformation was small; they retain coefficient one.
This endpoint can start the next lift.

Each nontrivial stratum surgery forces at least one negative step
in its lift.  Its negative morphism is nontrivial as well, by
smallness, matched boundary data, and the ample signs; lift a
negative curve to see the assertion.  Infinitely many surgeries
would concatenate to an infinite lower-dimensional dlt elementary
sequence across possibly varying bases, contrary to (iii) in
dimension $e$.  This proves disjointness for the chosen center.
There are finitely many centers, completing (i) in dimension $d$.

For (ii), take a projective SNC resolution of the lc pair with
boundary equal to the strict boundary plus the full reduced
exceptional divisor.  Its adjoint is the pullback of the original
adjoint plus an effective exceptional correction, by log canonicity;
the correction is supported in the floor.  Run relative elementary
steps over the original space.  At every stage the same identity
holds with the pushed-forward correction.  If the run were
infinite, (i) and the cycle count would leave a tail disjoint from
the floor on both sides.  The current adjoint could not have
negative degree on a contracted curve there, since its correction
is supported in the floor and its other term is a base pullback.
Thus the run ends.  Relative nefness and negativity kill its
effective exceptional correction.  The endpoint is crepant dlt
and globally strongly $\Q$-factorial by continuation.  Every
remaining exceptional prime came from an exceptional prime on
the resolution and still has coefficient one.  This proves (ii).

Finally, if $d\leq3$, an infinite dlt sequence would, by (i) and the
cycle count, have a small tail disjoint from the floor on both
sides.  Delete the floor.  The resulting pair is klt, using the
global strong condition and the dlt discrepancy characterization.
The ample signs are unchanged by disjointness on both sides.
This contradicts Theorem~\ref{bir:klt-termination}, proving (iii)
and completing the induction.
\end{proof}

\subsection{A reduced-boundary model with a nef klt interval}
\label{bir:reduced-section}

\begin{lemma}[Transport of canonical pseudo-effectivity]
\label{bir:pseudoeffective-transport}
Suppose a bimeromorphic transformation of normal globally strongly
$\Q$-factorial spaces extracts no divisors, and a Cartier power of
the canonical line on its source has a semipositive singular metric.
Then a Cartier power of the canonical line on its target has such
a metric as well.
\end{lemma}

\begin{proof}
Choose a common canonical power on the two spaces.  On their
isomorphic open, the actual canonical identification transports the
metric.  Nonextraction makes the complement of this open in the
target have codimension at least two.  In a local frame of the
target line we must extend a psh weight over that complement.

Here is the local upper bound needed for this Hartogs argument on a
normal germ.  Take a finite local projection to a ball.  Off the
branch locus and the image of the missing set, the maximum of the
weight on the finite fibers is psh.  It extends across the branch
locus while the fibers still avoid the missing set, by local upper
boundedness.  The bad image has codimension at least two, so psh
Hartogs extension on the smooth ball extends this maximum there as
well.  It bounds the original weight on a dense analytic complement.
The bound holds on its whole original domain by the disc test,
using discs generically outside the excluded analytic sets through
any tested point.  Normality and the psh Riemann extension theorem
on locally irreducible spaces now give the psh extension by upper
regularization.  These extensions respect the transition functions
of the actual line.  They define the required semipositive metric.
For a related actual-line statement under rational singularities,
see \cite[Lemma~3.6]{HLLNonvanishing}.
\end{proof}

\begin{proposition}[Reduced-boundary model]
\label{bir:reduced-model}
Assume Assumption~\ref{hyp:mmp}.  Let $(T_0,G_0)$ be an ordinary
rational dlt fourfold pair with reduced boundary, on a normal
irreducible globally strongly $\Q$-factorial compact K\"ahler space.
Assume that $K_{T_0}$ has a semipositive singular metric on a Cartier
power.  The boundary may be empty.  There is a nonextracting
bimeromorphic map $T_0\dashrightarrow T$ such that:
\begin{enumerate}[label=\textup{(\roman*)}]
\item $T$ is normal, compact K\"ahler, globally strongly
$\Q$-factorial, and klt for zero boundary; $K_T$ remains
pseudo-effective with such a semipositive metric.
\item The pushforward $G$ is reduced, $(T,G)$ is lc, and the
actual rational line $A=K_T+G$ is analytically nef.  There is a
rational $t_0\in[0,1)$ such that $(T,tG)$ is klt and
$K_T+tG$ is analytically nef for every rational
$t\in[t_0,1)$.
\item There is a projective crepant dlt modification
\begin{equation}
 h:(V,D)\longrightarrow(T,G),\qquad
 J:=K_V+D=h^*A,
 \label{bir:reduced-adjoint}
\end{equation}
where $V$ is normal, compact K\"ahler, and globally strongly
$\Q$-factorial, and $D$ is reduced.  The line $J$ is analytically
nef.  If $G\ne0$, then $D\ne0$.
\item On a simultaneous smooth compact K\"ahler projective
resolution $\mu:M\to V$ of the displayed models and boundaries,
the line $\mu^*J$ has a semipositive metric with minimal
singularities and zero Lelong numbers everywhere.
\end{enumerate}
The construction uses the fourfold MMP assumption only for ordinary
effective klt pairs with pseudo-effective adjoint.
\end{proposition}

\begin{proof}
We give the two phases separately.  In the first phase the current
pair $(T_i,G_i)$ stays dlt with reduced boundary, and we make a step
whenever $K_{T_i}+G_i$ has a negative extremal ray of the cone
$\NA(T_i)$ specified in Assumption~\ref{hyp:mmp}.
The underlying zero-boundary pair is klt by the dlt floor
perturbation argument, and canonical pseudo-effectivity is retained
by Lemma~\ref{bir:pseudoeffective-transport}.

For a chosen negative ray, take a rational $t<1$ sufficiently close
to one that it stays negative for $K_{T_i}+tG_i$.  This is an
effective klt pair, and its adjoint is pseudo-effective because
$K_{T_i}$ is.  Assumption~\ref{hyp:mmp} supplies the projective
bimeromorphic ray step with its prescribed face, negative-side
Bott--Chern rank, global strong condition, and K\"ahler category.
The full adjoint has negative degree on this ray, hence the negative
relative ample sign.  Its one-ray degree-zero adjustment and
Lemma~\ref{bir:integral-descent}, applied to the klt contraction,
give the positive relative ample sign on a flip.  Thus it is an
elementary dlt step for the full adjoint, and
Lemma~\ref{bir:comparison} keeps the transformed pair dlt.

This first phase is finite.  Otherwise the cycle count and
Theorem~\ref{bir:special-termination}(i) leave a small tail disjoint
from the boundary on both sides.  It is then a genuine
zero-boundary $K$-negative klt program.  The same rays are
$K$-negative, as tested by their actual contracted curves; the
contraction faces and negative-side Bott--Chern conditions are the
ones already supplied.  Both relative canonical signs are unchanged
because the floor is disjoint from all fibers involved in the
surgery.  Moreover these are the actual canonical flips stipulated
by the assumption: under a small map, direct images of common
Cartier canonical powers agree as reflexive sheaves, and the
positive canonical powers are relatively ample.  They give the
required relative canonical algebra and tautological line.  The
supporting-class existence for these zero-boundary negative rays
is also supplied by the same assumption.  This infinite tail would
therefore be one of the programs excluded by its arbitrary
termination clause, starting from its first zero-boundary klt
model.  No dlt fourfold termination assertion has been used.

At the end of phase one write $A_i=K_i+G_i$.  It has no negative
extremal ray and, in particular, is nonnegative on curves.  We
justify this conclusion without asserting an absolute curve
criterion for nefness.  The normalized K\"ahler-mass slice of the
positive closed current cone defining
$\NA(T_i)$ is compact for smooth Bott--Chern pairings.
Indeed positive closed currents of unit mass have locally bounded
order-zero coefficients by positivity and the K\"ahler mass; weak
limits remain positive and closed.  Their pairing image is compact
in the product of the pairing coordinates, or in the
finite-dimensional realization when used.  Its positive cone with
zero is closed, since the mass coordinate and normalized
coordinates control every convergent sequence.  Thus this is the
full cone in the definition.  A negative value of $A_i$ would give
a negative minimum face of the compact convex slice and hence an
extreme point, contradicting the absence of a negative extremal
ray.  Only curve nonnegativity is needed at this stage.

In the second phase we make $A_i$-trivial steps of a single
zero-boundary $K$ program.  Maintain the properties that $A_i$ is
nonnegative on curves, $(T_i,G_i)$ is lc, $(T_i,0)$ is klt, and
$K_i$ is pseudo-effective.  If some rational $t_0\in[0,1)$ already
has $K_i+t_0G_i$ analytically nef, every rational $t\in[t_0,1)$
does as well.  Indeed these pairs are klt by interpolation between
the zero-boundary klt pair and $(T_i,G_i)$, and their adjoints are
pseudo-effective.  Their degrees are nonnegative by convexity
between $K_i+t_0G_i$ and the curve-nonnegative $A_i$.  A non-nef
one would, by Assumption~\ref{hyp:mmp}, have a negative ray with a
projective contraction and an actual negative contracted curve, a
contradiction.  Closedness of the analytic nef cone then makes
$A_i$ analytically nef as $t$ tends to one.

Suppose instead that every rational $K_i+tG_i$, $0\leq t<1$, is
non-nef.  Choose current positive integers $m,r$ with $mA_i$ and
$rK_i$ Cartier.  Choose a rational $t$ so close to one that
\begin{equation}
 \frac{m(1-t)}{t}<\frac{1}{r(\dim T_i+1)}.
 \label{bir:near-one-choice}
\end{equation}
The klt assumption gives a negative extremal ray $R$ for
\[
 K_i+tG_i=tA_i+(1-t)K_i.
\]
Its contraction supplies an actual curve class.  Since $A_i$ is
nonnegative on curves, $R$ is $K_i$-negative.  Choose the prescribed
contraction and step for this same ray now from the zero-boundary
application of Assumption~\ref{hyp:mmp}.

We claim $A_i\cdot R=0$.  If this degree were positive, the Cartier
line $mA_i$ would be relatively ample on that contraction, since
its fiber-curve degrees are a positive multiple of those of $-K_i$.
Work over a Stein neighborhood with a property-(P) compactum
containing a point with nontrivial fiber.  The local projective
analytic rationality theorem for the ample Cartier line $mA_i$
and the klt canonical adjoint gives denominator at most
$r(\dim T_i+1)$ for its positive finite relative nef threshold;
see \cite[Theorem~4.3.1 and Remark~4.3.4]{FujinoAnalyticCone}.
All contracted curve degrees are proportional, so for a curve $C$
of the ray that threshold is exactly
\begin{equation}
 \frac{mA_i\cdot C}{-K_i\cdot C}
 \geq\frac{1}{r(\dim T_i+1)}.
 \label{bir:rationality-bound}
\end{equation}
The actual line data and canonical representatives needed for the
theorem are available locally over this bimeromorphic Stein base,
as in Section~\ref{bir:finite-generation}.  On the other hand,
negativity for the test boundary gives
\[
 \frac{mA_i\cdot C}{-K_i\cdot C}
 <\frac{m(1-t)}{t},
\]
contradicting~\eqref{bir:near-one-choice}.  This proves the claim
using only current-stage indices, with no uniform index assumption
on the fourfold sequence.

By Lemma~\ref{bir:integral-descent}, a Cartier multiple of $A_i$
descends to an actual line on the contraction base.  The pushed
adjoint on a divisorial step, or the transformed adjoint on a
small positive model, is the corresponding pullback by
codimension-one identification and reflexivity.  On a common
projective model any remaining exceptional comparison is
relatively numerically zero, and hence zero by negativity.
Thus the step is crepant for the actual adjoint $A_i$, and the
pair stays lc.  Curve nonnegativity on the next model follows by
lifting each curve to a common model and using the equality of
pullbacks.  The zero-boundary klt and canonical pseudo-effective
hypotheses persist by Assumption~\ref{hyp:mmp} and
Lemma~\ref{bir:pseudoeffective-transport}.

Repeat this procedure whenever no rational nef parameter exists.
Every step is a prescribed $K_i$-negative step of the same
zero-boundary program.  Arbitrary termination in
Assumption~\ref{hyp:mmp} excludes indefinite repetition.  We
therefore reach the asserted $T,G,A$ and nef klt interval.  Neither
phase extracts divisors, so their composite is nonextracting.

Apply Theorem~\ref{bir:special-termination}(ii) to $(T,G)$ to obtain
\eqref{bir:reduced-adjoint}.  The boundary $D$ is the strict
transform of $G$ plus coefficient-one exceptional primes, hence
reduced.  If $G$ is nonzero its strict transform is nonzero.
Analytic nefness pulls back, giving that of $J$.

Finally take the stated simultaneous resolution and write
$\rho=h\circ\mu:M\to T$.  For every rational $t$ in the nef klt
interval, Lemma~\ref{met:zero-lelong} applied to
$(T,tG)$ gives a semipositive metric with minimal singularities
and zero Lelong numbers on $\rho^*(K_T+tG)$.  Add the metric of
the effective rational divisor $(1-t)\rho^*G$.  This is a
semipositive metric on the fixed actual line
\begin{equation}
 \rho^*A=\rho^*(K_T+tG)+(1-t)\rho^*G=\mu^*J.
 \label{bir:metric-interval}
\end{equation}
A metric with minimal singularities on that line is no more
singular than each of these metrics, up to an additive constant
in weights.  At every $x\in M$ its Lelong number is therefore at
most $(1-t)\nu([\rho^*G],x)$.  The latter is finite and tends to
zero as rational $t$ tends to one.  All Lelong numbers of the
minimal metric vanish.  This proves (iv) and the proposition.
\end{proof}
